%% =====================================================================
%% Canovaccio della discussione NDS (prof. Baggio), 26/11/2026, ore 9:30
%% Solo il discorso: per ogni slide le frasi da dire, nell'ordine della slide.
%% =====================================================================
%% COME SI AGGIORNA
%% 1. Ogni slide e' un blocco
%%       \begin{slide}{<label del frame>}{<titolo>}{<tempo>}
%%         \begin{discorso}
%%           \item prima frase
%%           \item seconda frase
%%         \end{discorso}
%%       \end{slide}
%%    La <label> e' quella del frame in main.tex, es. [label=m-bridge].
%%    Frame a due passi: \begin{slidedue}{...}{...}{...}; \clic segna il passo 2.
%%    Pagine senza label (titolo, outline, Thank you): \begin{slidepag}{<pagina>}{...}{...}.
%% 2. Miniature, numeri di slide e di pagina NON si scrivono a mano: vengono
%%    da ../main.pdf e ../main.aux. Se cambi la presentazione, ricompila prima
%%    main.tex e poi questo file (pdflatex canovaccio.tex, due volte).
%% 3. Nel testo: \fr{label} = numero della slide; \pg{label} = pagina del PDF.
%% 4. Domande probabili, termini da evitare e cose nostre: file domande.tex.
%% =====================================================================
\documentclass[a4paper,10pt]{article}
\usepackage[T1]{fontenc}
% babel in italiano: con TeX Live completo usa italian.ldf, altrimenti il file .ini di babel
\IfFileExists{italian.ldf}{\usepackage[italian]{babel}}%
  {\usepackage{babel}\babelprovide[import,main]{italian}}
\usepackage{lmodern}
\usepackage{microtype}
\usepackage[a4paper,margin=1.5cm,top=1.4cm,bottom=1.7cm]{geometry}
\usepackage{amsmath,amssymb}
\usepackage{graphicx}
\usepackage{xcolor}
\usepackage{enumitem}
\usepackage{refcount}
\usepackage{xr-hyper}
\usepackage[hidelinks]{hyperref}
\externaldocument[P-]{../main}  % legge ../main.aux
\pdfsuppresswarningpagegroup=1

% ---------------------------------------------------------------------
% Riferimenti alla presentazione
\newcommand{\presentazione}{../main.pdf}
\newcommand{\pg}[1]{\getpagerefnumber{P-#1}}  % pagina del PDF
\newcommand{\fr}[1]{\getrefnumber{P-#1}}      % numero della slide
\newcommand{\uno}{\mathbf{1}}

% ---------------------------------------------------------------------
% Impaginazione
\setlength{\parindent}{0pt}
\setlength{\parskip}{2pt}
\raggedbottom
\definecolor{unipd}{RGB}{155,0,20}

\newcommand{\mini}[1]{{\setlength{\fboxsep}{0pt}%
  \fbox{\includegraphics[page=#1,width=\dimexpr\linewidth-2\fboxrule\relax]{\presentazione}}}}
\newcommand{\testata}[4]{% numero, pagine, titolo, tempo
  {\color{unipd}\bfseries Slide #1 \textperiodcentered\ #3}\hfill{\small\itshape #4}\par
  {\scriptsize\color{gray} pagina #2 del PDF}\par\smallskip}

% Blocco di una slide: miniatura a sinistra, discorso a destra
\newenvironment{slide}[3]{%
  \par\addvspace{12pt}\noindent
  \begin{minipage}[t]{0.36\textwidth}\vspace{0pt}\mini{\pg{#1}}\end{minipage}\hfill
  \begin{minipage}[t]{0.61\textwidth}\vspace{0pt}%
  \testata{\fr{#1}}{\pg{#1}}{#2}{#3}}%
 {\end{minipage}\par}
% Frame a due passi: le due miniature una sotto l'altra
\newenvironment{slidedue}[3]{%
  \par\addvspace{12pt}\noindent
  \begin{minipage}[t]{0.36\textwidth}\vspace{0pt}\mini{\pg{#1<1>}}\par\smallskip
    \mini{\pg{#1<2>}}\end{minipage}\hfill
  \begin{minipage}[t]{0.61\textwidth}\vspace{0pt}%
  \testata{\fr{#1}}{\pg{#1<1>}--\pg{#1<2>}}{#2}{#3}}%
 {\end{minipage}\par}
% Pagine senza label: si indica la pagina del PDF
\newenvironment{slidepag}[3]{%
  \par\addvspace{12pt}\noindent
  \begin{minipage}[t]{0.36\textwidth}\vspace{0pt}\mini{#1}\end{minipage}\hfill
  \begin{minipage}[t]{0.61\textwidth}\vspace{0pt}%
  {\color{unipd}\bfseries #2}\hfill{\small\itshape #3}\par
  {\scriptsize\color{gray} pagina #1 del PDF}\par\smallskip}%
 {\end{minipage}\par}

% Il discorso: una frase (o due) per punto
\newenvironment{discorso}{\begin{itemize}[nosep,leftmargin=1em,label={\color{gray}\textbullet},itemsep=2pt]}{\end{itemize}}
\newcommand{\clic}{{\color{unipd}\scriptsize[passo 2]}\ }
\newcommand{\parte}[2]{\section*{#1\hfill{\normalsize\itshape #2}}}

\begin{document}

\begin{center}
  {\LARGE\bfseries Canovaccio della discussione}\\[3pt]
  {\large Stubbornness and Social Power in Opinion Dynamics}\\[2pt]
  Network Dynamical Systems, prof.\ Giacomo Baggio \textperiodcentered\ giovedì 26 novembre 2026, ore 9:30\\[2pt]
  {\small\color{gray} Versione del \today}
\end{center}

% =====================================================================
\parte{Parte 1 — Introduzione}{3:00}

\begin{slidepag}{1}{Titolo}{0:15}
  \begin{discorso}
    \item Buongiorno. Presento il mio progetto su stubbornness e social power nella dinamica delle opinioni.
    \item Studio due modelli, il modello di French–Harary–DeGroot e il modello di Friedkin–Johnsen: come cambia il loro comportamento quando gli agenti resistono agli altri, e quanto pesa ciascun agente nelle opinioni finali.
  \end{discorso}
\end{slidepag}

\begin{slidepag}{2}{Slide 2 \textperiodcentered\ Outline}{0:15}
  \begin{discorso}
    \item La presentazione ha quattro parti: introduzione, modello FDG, modello FJ, conclusioni.
    \item Alcune slide hanno dei pulsanti in basso a destra: aprono le slide di backup, con definizioni e derivazioni. Le uso solo se vuole entrare nel dettaglio.
  \end{discorso}
\end{slidepag}

\begin{slide}{m-question}{Stubbornness}{1:00}
  \begin{discorso}
    \item Una rete sociale è un digrafo pesato: gli agenti sono i nodi, gli archi portano l'influenza, e ogni agente aggiorna la sua opinione facendo la media delle opinioni che ascolta.
    \item Un agente è stubborn se la sua opinione resta invariata indipendentemente dalle opinioni degli altri: è la definizione del tutorial di Proskurnikov e Tempo.
    \item È il caso estremo di due modi graduati di resistere agli altri: dare peso alla propria opinione attuale, oppure restare attaccati a un'opinione fissa propria, un pregiudizio.
    \item Il primo è il modo del modello FDG, il secondo del modello FJ: li introduco tra poco.
    \item A entrambi i modelli, sulle stesse reti, faccio tre domande. Le opinioni convergono? Convergono alla stessa opinione per tutti, il consenso, o a opinioni diverse, il disaccordo? E quanto pesa l'opinione iniziale di ciascun agente in quelle finali?
    \item Prima di tutto, da dove vengono i due modelli.
  \end{discorso}
\end{slide}

\begin{slide}{m-models}{Models of opinion dynamics}{0:40}
  \begin{discorso}
    \item Il quadro storico è quello del tutorial di Proskurnikov e Tempo, la fonte principale del lavoro, insieme alle slide del corso e al libro di Bullo.
    \item A sinistra i modelli che fanno solo medie: in tempo discreto il modello di French, Harary e DeGroot; in tempo continuo quello di Abelson.
    \item Per i gruppi reali il consenso è una previsione troppo forte: spesso restano divisi.
    \item Per questo il modello di Friedkin e Johnsen, negli anni Novanta, aggiunge alle medie opinioni esterne costanti; è stato validato sperimentalmente su gruppi piccoli e medi. In tempo continuo il corrispondente è il modello di Taylor.
    \item In questo lavoro tratto il modello FDG, visto nel corso, e il modello FJ, che è nuovo; Abelson e Taylor restano solo nella panoramica.
    \item Richiamo solo le nozioni del corso che mi servono.
  \end{discorso}
\end{slide}

\begin{slide}{m-prelim}{Preliminaries}{0:50}
  \begin{discorso}
    \item Sono nozioni del corso, quindi solo i punti che userò.
    \item Un arco da $i$ a $j$, con peso $a_{ij}$, vuol dire che l'agente $i$ è influenzato dall'agente $j$: è la convenzione del corso; il tutorial disegna gli archi al contrario.
    \item Un digrafo fortemente connesso è periodico se un intero maggiore di uno divide la lunghezza di ogni ciclo semplice; un self-loop lo rende aperiodico.
    \item Nel condensation digraph ogni componente fortemente connessa diventa un nodo; un sink è un gruppo influenzato solo da sé stesso. Saranno i sink a decidere l'esito.
    \item Per le matrici: semi-convergente vuol dire che il limite delle potenze esiste; convergente, o Schur stabile, che il limite è zero, cioè raggio spettrale minore di uno.
    \item Infine: irriducibile equivale a fortemente connesso, primitiva a fortemente connesso e aperiodico.
    \item Ora il primo modello.
  \end{discorso}
\end{slide}

% =====================================================================
\parte{Parte 2 — Il modello FDG}{7:00}

\begin{slide}{m-fdg-model}{The FDG model}{0:50}
  \begin{discorso}
    \item Nel modello FDG ogni agente sostituisce la sua opinione con una media pesata delle opinioni che ascolta.
    \item In forma vettoriale $x(k+1)=Ax(k)$, con $A$ la matrice di adiacenza pesata, row-stochastic: è un sistema di media a tempo discreto.
    \item Iterando, $x(k)=A^kx(0)$: il comportamento nel lungo periodo dipende dalle potenze di $A$. Una matrice row-stochastic ha sempre l'autovalore 1, e raggio spettrale 1.
    \item Il modello è convergente se, per ogni condizione iniziale, il limite esiste; raggiunge un consenso se, per ogni condizione iniziale, le opinioni finali sono tutte uguali.
    \item Gli esiti possibili sono quindi tre: nessuna convergenza, cioè $A$ non semi-convergente; consenso; e disaccordo, cioè il modello converge ma per qualche condizione iniziale le opinioni finali non sono tutte uguali.
    \item Quale dei tre si ha si legge in tre modi equivalenti.
  \end{discorso}
\end{slide}

\begin{slide}{m-fdg-rosetta}{Which outcome: the Rosetta stone}{0:40}
  \begin{discorso}
    \item L'esito si legge sullo spettro di $A$, sulla matrice o sul digrafo: la Rosetta stone del corso collega i tre punti di vista.
    \item Sul digrafo conta solo il condensation digraph: quanti sink ha e se sono periodici.
    \item Un solo sink, aperiodico, dà il consenso: casi (i) e (iii). Almeno due sink, tutti aperiodici, danno il disaccordo: caso (iv). Un sink periodico impedisce la convergenza: casi (ii) e (v).
    \item Mi servirà soprattutto il Corollario 1.12: se il digrafo è fortemente connesso e ha almeno un self-loop, $A$ è primitiva.
    \item Il lavoro si concentra sul consenso, perché è il caso in cui il peso di ciascun agente si può misurare. Al disaccordo dedico la slide dopo.
  \end{discorso}
\end{slide}

\begin{slide}{m-fdg-consensus}{Consensus and social power}{0:50}
  \begin{discorso}
    \item Il teorema del consenso del corso dà tre condizioni equivalenti.
    \item Sullo spettro: l'autovalore 1 è semplice e gli altri hanno modulo minore di uno. Sulla matrice: le potenze di $A$ convergono a $\uno w^\top$. Sul grafo: esiste un nodo globally reachable, e il sottografo dei nodi globally reachable è aperiodico.
    \item In questo caso $w$ è l'autovettore sinistro dominante di $A$, con $w_i$ positivo se e solo se il nodo $i$ è globally reachable, e il valore di consenso è $w^\top x(0)$.
    \item $w_i$ è la social power dell'agente $i$: il peso della sua opinione iniziale nel valore di consenso.
    \item Richiede il consenso, perché solo allora le righe del limite di $A^k$ sono tutte uguali a $w^\top$.
    \item Se $A$ è doubly-stochastic e vale la condizione sul grafo, $w$ è uniforme e il consenso è la media delle opinioni iniziali.
    \item Nei termini del tutorial: il modello converge se e solo se $A$ è regular, e raggiunge il consenso se e solo se $A$ è fully regular.
    \item E se il consenso non c'è?
  \end{discorso}
\end{slide}

\begin{slide}{m-fdg-disagreement}{Disagreement and stubbornness}{0:40}
  \begin{discorso}
    \item Se il condensation digraph ha almeno due sink, il teorema del disaccordo dà tre condizioni equivalenti: l'autovalore 1 è semi-semplice, con molteplicità pari al numero di sink; $A$ è semi-convergente; ogni sink è aperiodico.
    \item Allora ogni sink raggiunge un consenso al suo interno, e ogni agente fuori dai sink converge a una combinazione convessa dei valori dei sink.
    \item Le righe del limite di $A^k$ sono diverse: nessun vettore fa la parte di $w$, e la social power non è definita.
    \item A destra la definizione di stubborn del tutorial. Nel modello FDG l'agente $i$ è stubborn se e solo se $a_{ii}=1$: è un sink fatto di un solo nodo, aperiodico grazie al self-loop.
    \item Se tutti gli agenti sono influenzati da lui, direttamente o attraverso altri, è l'unico nodo globally reachable, e la sua social power è 1.
    \item Ora le due reti su cui confronto i due modelli.
  \end{discorso}
\end{slide}

\begin{slide}{m-networks}{Two networks, one agent}{0:40}
  \begin{discorso}
    \item Uso due reti di sei agenti: un ciclo diretto, in cui ogni agente ascolta solo il successivo, e un anello bidirezionale, in cui ogni agente ascolta i due vicini con peso un mezzo.
    \item Per $t=0$ entrambe sono fortemente connesse e periodiche, con periodo 6 il ciclo e 2 l'anello, e doubly-stochastic: l'autovettore sinistro dell'autovalore 1 è uniforme, nessun agente è favorito.
    \item Cambio solo la riga dell'agente 1: $t$ è il peso che l'agente 1 dà alla propria opinione attuale. Si va dalle reti originali, $t=0$, all'agente 1 stubborn, $t=1$.
    \item La tabella in basso: per $t$ fra 0 e 1 le reti sono fortemente connesse e aperiodiche, e non più doubly-stochastic; per $t=1$ non sono più fortemente connesse, perché l'agente 1 non ascolta nessun altro.
    \item La condizione iniziale è sempre la stessa: opinioni da 0 a 1, con passo 0.2.
    \item Vediamo cosa succede alle opinioni.
  \end{discorso}
\end{slide}

\begin{slide}{m-fdg-conv}{One self-weight brings agreement}{0:50}
  \begin{discorso}
    \item Le colonne sono $t=0$, un mezzo e 1; sopra il ciclo, sotto l'anello.
    \item Per $t=0$ le reti sono periodiche e le opinioni oscillano per sempre: nessuna convergenza, caso (ii). Sullo spettro ci sono autovalori diversi da 1 sul cerchio unitario.
    \item Basta un self-weight $t$ positivo, di un solo agente, e il gruppo raggiunge il consenso.
    \item Per $t$ fra 0 e 1 la rete è fortemente connessa con un self-loop, quindi $A(t)$ è primitiva per il Corollario 1.12: caso (i).
    \item Per $t=1$ l'agente 1 è stubborn ed è l'unico nodo globally reachable: tutti finiscono sulla sua opinione iniziale, zero. Caso (iii).
    \item Per ogni $t$ positivo gli autovalori diversi da 1 sono tutti strettamente dentro il disco unitario.
    \item Questo è il primo punto del lavoro: la testardaggine ribalta la dinamica. Un agente che dà peso alla propria opinione attuale, senza essere stubborn, basta perché su una rete fortemente connessa tutto il gruppo si metta d'accordo.
    \item Con il consenso possiamo misurare il peso di ciascun agente.
  \end{discorso}
\end{slide}

\begin{slide}{m-fdg-sp}{Social power grows gradually}{0:45}
  \begin{discorso}
    \item Per $t$ in $(0,1]$ c'è il consenso, e la social power ha la forma chiusa a destra, la stessa sulle due reti.
    \item In alto: la social power dell'agente 1 cresce gradualmente con il suo self-weight, da $1/6$, che è il limite per $t$ che tende a 0, fino a 1 per $t=1$. A $t=\tfrac12$ vale $2/7$, e ogni altro agente $1/7$.
    \item In basso il valore di consenso, $3(1-t)/(6-5t)$: da 0.5 nel limite per $t$ che tende a 0, fino a $x_1(0)=0$; a $t=\tfrac12$ vale $3/7$.
    \item I valori sono gli stessi sulle due reti: la rete decide come le opinioni convergono, non dove.
    \item In termini sociali: l'inerzia guadagna influenza poco alla volta, e per $t$ fra 0 e 1 ogni opinione iniziale entra ancora nel consenso.
    \item Perché la stessa social power su due reti così diverse?
  \end{discorso}
\end{slide}

\begin{slide}{m-fdg-why}{Why the same social power}{0:45}
  \begin{discorso}
    \item La social power risolve $w^\top A(t)=w^\top$, che si legge come un bilancio: ogni agente passa la sua social power agli agenti che ascolta, in proporzione ai pesi, e $w_j$ è quello che l'agente $j$ raccoglie.
    \item $A(0)$ è doubly-stochastic: ogni agente è ascoltato con peso totale 1, quindi il bilancio torna quando ogni agente passa la stessa quantità.
    \item Gli agenti da 2 a 6 non hanno self-loop e passano tutto il loro $b$; l'agente 1 trattiene la quota $t$ e passa $(1-t)w_1$.
    \item Quindi $(1-t)w_1=b$: l'agente 1 pesa $1/(1-t)$ volte ogni altro agente, $2/7$ contro $1/7$ a $t=\tfrac12$.
    \item Conta solo quanto trattiene l'agente 1, non chi ascolta: per questo $w(t)$ è la stessa sulle due reti.
    \item A destra la verifica del report: sostituendo $w=(w_1,b,\dots,b)$ si ritrova $b=(1-t)w_1$.
    \item Per $t$ che tende a 0, $w$ tende al vettore uniforme; ma a $t=0$ non c'è consenso, e quel vettore è solo l'autovettore sinistro di $A(0)$.
    \item Finora al massimo un agente stubborn. Vediamo cosa succede con due.
  \end{discorso}
\end{slide}

\begin{slide}{m-bridge}{Two stubborn agents}{1:00}
  \begin{discorso}
    \item Come nella Figura 2 con $t=1$, l'agente 1 è stubborn; ora sull'anello lo è anche l'agente 4, quello opposto.
    \item Il condensation digraph ha due sink, $\{1\}$ e $\{4\}$, aperiodici per i self-loop.
    \item Per il Teorema 2.4 il gruppo converge a un disaccordo: gli agenti 2, 3, 5 e 6 finiscono su combinazioni convesse delle opinioni iniziali di 1 e di 4.
    \item Le righe del limite di $A^k$ sono diverse, quindi la social power non è definita.
    \item Qui mi fermo con il modello FDG, per due motivi: è già visto nel corso, e la social power, che è ciò che mi interessa, richiede il consenso.
    \item Però lascio aperta una domanda. Gli agenti 1 e 4 non cambiano mai opinione: sono membri del gruppo, o sono opinioni fisse, esterne, che gli altri ascoltano?
    \item E in generale: un agente la cui opinione resta invariata indipendentemente dalle opinioni degli altri è un membro del gruppo, o un'opinione fissa esterna che gli altri ascoltano?
    \item Ci torno più avanti. Ora il secondo modello, il modello di Friedkin–Johnsen.
  \end{discorso}
\end{slide}

% =====================================================================
\parte{Parte 3 — Il modello FJ, teoria}{6:00}

\begin{slide}{m-fj-model}{The FJ model}{1:00}
  \begin{discorso}
    \item Nel modello FJ ogni agente pesa con $\lambda_i$ la media delle opinioni che ascolta, e con il complemento $1-\lambda_i$ il suo pregiudizio $u_i$.
    \item $W$ è row-stochastic, la matrice delle influenze sociali: ha il ruolo che $A$ aveva nel FDG.
    \item $\lambda_i$ è la suscettibilità dell'agente $i$ all'influenza sociale; $u$ è il vettore costante dei pregiudizi, opinioni fisse che gli agenti ascoltano; $1-\lambda_i$ è l'attaccamento al proprio pregiudizio.
    \item Se $\lambda_i=1$ l'agente segue la regola del FDG; con $\Lambda=I$ il modello è proprio il FDG con $A=W$, quindi da qui in poi $\Lambda\neq I$.
    \item Se $\lambda_i<1$ l'agente è prejudiced. Se $\lambda_i=0$ l'agente vale $u_i$ da $k=1$ in poi, ma è stubborn solo se in più $u_i=x_i(0)$.
    \item Nel modello FJ originale $u=x(0)$ e $1-\lambda_i=w_{ii}$; come nel tutorial, qui non lo impongo.
    \item La differenza con il FDG sta nell'oggetto: un self-weight lega un agente alla sua opinione attuale, che si muove ancora verso gli altri; un pregiudizio lo lega a un'opinione fissa, che non si muove mai.
    \item Torniamo alla prima domanda: quando converge?
  \end{discorso}
\end{slide}

\begin{slide}{m-fj-conv}{Convergence in the FJ model}{1:00}
  \begin{discorso}
    \item Il modello FJ converge se il limite esiste per ogni condizione iniziale, ed è asintoticamente stabile se $\Lambda W$ è convergente, cioè ha raggio spettrale minore di uno.
    \item La nozione chiave è quella di agente P-dependent: un agente prejudiced, oppure influenzato, direttamente o attraverso altri, da un agente prejudiced. Gli altri sono P-independent.
    \item Gli agenti P-independent non ascoltano nessun P-dependent: formano un modello FDG a sé, che nessun pregiudizio raggiunge.
    \item Il Teorema 3.2 dice che il blocco dei P-dependent è sempre asintoticamente stabile, e che il modello converge se e solo se sono tutti P-dependent oppure $W^{22}$ è regular. Le opinioni finali sono combinazioni convesse dei pregiudizi e delle opinioni finali dei P-independent.
    \item Il caso che userò di più è il Corollario 3.3: il modello è asintoticamente stabile se e solo se tutti gli agenti sono P-dependent. Per esempio quando la rete è fortemente connessa e almeno un agente è prejudiced. E non serve nessuna aperiodicità.
    \item L'intuizione: lo scarto dall'equilibrio evolve come $e(k+1)=\Lambda W\,e(k)$, senza $u$. Un agente prejudiced ne perde a ogni passo una frazione $1-\lambda_i$, quindi gli scarti si spengono solo se ogni agente è influenzato da un agente prejudiced.
    \item Seconda domanda: dove convergono le opinioni?
  \end{discorso}
\end{slide}

\begin{slide}{m-fj-final}{Final opinions and the augmented graph}{1:00}
  \begin{discorso}
    \item Con tutti gli agenti P-dependent le opinioni finali sono $x(\infty)=Vu$.
    \item Tre conseguenze. Le opinioni iniziali sono dimenticate: $x(0)$ non compare. Contano solo i pregiudizi degli agenti prejudiced: se $\lambda_j=1$, la colonna $j$ di $V$ è nulla.
    \item E, per un dato $u$, il gruppo raggiunge un consenso se e solo se tutti gli agenti prejudiced hanno lo stesso pregiudizio.
    \item Il tutorial dà anche una seconda lettura, il modello aumentato: a ogni agente $i$ si affianca un agente virtuale stubborn $i'$, che vale sempre $u_i$.
    \item La matrice $\bar A$ è row-stochastic, quindi è un modello FDG su $2n$ agenti. Nella figura l'agente 1 è l'unico prejudiced: a sinistra il modello FJ, a destra il grafo aumentato con l'agente virtuale $1'$.
    \item Ogni agente virtuale è un sink aperiodico; con tutti gli agenti P-dependent sono gli unici sink. Letto con la Rosetta stone, è il caso (iv): il disaccordo del Teorema 2.4, con limite proprio $Vu$.
    \item In termini sociali: un pregiudizio si comporta come un agente in più che non cambia mai opinione, messo accanto all'agente prejudiced.
    \item Terza domanda: quanto pesa ciascun agente?
  \end{discorso}
\end{slide}

\begin{slide}{m-fj-centrality}{Influence centrality}{0:45}
  \begin{discorso}
    \item La social power richiede il consenso, che nel modello FJ non è una proprietà del modello: dipende da $u$. E il modello aumentato ha almeno $n$ sink, quindi per $n\ge2$ non va mai al consenso.
    \item Come nel tutorial, allora, prendo il modello FJ asintoticamente stabile con $u=x(0)$, l'ipotesi tradizionale della teoria FJ. Così $x(\infty)=Vx(0)$.
    \item La influence centrality $c_i$ è il peso medio dell'opinione iniziale dell'agente $i$ nelle opinioni finali del gruppo: $c=n^{-1}V^\top\uno$, e la media delle opinioni finali è $c^\top x(0)$.
    \item $V$ è stocastica, quindi $c$ è non negativo e somma a uno, come $w$.
    \item Con $\Lambda=\alpha I$: per $\alpha=0$ si ha $V=I$ e $c$ uniforme.
    \item Il Lemma 3.5 collega le due misure: se $W$ è fully regular, $c_\alpha$ tende alla social power del modello FDG con $A=W$ quando $\alpha$ tende a 1.
    \item Torno ora all'agente stubborn, visto nei due modelli.
  \end{discorso}
\end{slide}

\begin{slidedue}{m-state-input}{Stubborn agent: state or input}{1:15}
  \begin{discorso}
    \item Il tutorial mostra che un modello FDG con agenti stubborn si riscrive come un modello FJ: $\lambda_i=0$ e $u_i=x_i(0)$ per gli agenti stubborn, $\lambda_i=1$ per gli altri.
    \item A sinistra il FDG con $A(1)$, a destra il FJ con $\lambda_1=0$: le traiettorie sono le stesse.
    \item Ma la stabilità no. Il modello FDG converge ma non è mai asintoticamente stabile, perché 1 è sempre un autovalore di $A$; il modello FJ invece è asintoticamente stabile, per il Corollario 3.3.
    \item \clic Per vederlo perturbo l'opinione iniziale dell'agente 1: $x_1(0)=0.5$, con $u$ invariato.
    \item Nel FDG l'agente 1 tiene 0.5 e il gruppo lo segue: la perturbazione non viene mai dimenticata. Nel FJ l'agente 1 torna su $u_1=0$ al primo passo, e il limite è lo stesso, $Vu$.
    \item Sulle matrici: la riga dell'agente 1 è $e_1^\top$ in $A(1)$ ed è zero in $\Lambda W$, quindi l'autovalore 1 di $A(1)$ diventa 0 in $\Lambda W$.
    \item In termini sociali: nel FDG la stubbornness è una memoria nello stato, mai dimenticata; nel FJ è un ingresso esterno, e lo stato dimentica $x(0)$.
    \item Adesso posso rispondere alla domanda lasciata aperta.
  \end{discorso}
\end{slidedue}

\begin{slide}{m-back-question}{Back to the question}{1:00}
  \begin{discorso}
    \item Torno all'anello con gli agenti 1 e 4 stubborn.
    \item Nel report c'è una seconda riscrittura, nella Figura 8, lì con il solo agente 1: si tolgono gli agenti stubborn, e le loro opinioni diventano i pregiudizi degli agenti che li ascoltavano.
    \item Per esempio l'agente 2 aggiorna come $x_2(k+1)=\tfrac12x_3(k)+\tfrac12x_1(0)$: è il modello FJ con $\lambda_2=\tfrac12$ e $u_2=x_1(0)$.
    \item Allo stesso modo $u_6=x_1(0)$ e $u_3=u_5=x_4(0)$, tutti con $\lambda_i=\tfrac12$. Le traiettorie sono le stesse.
    \item Quindi la risposta è: valgono entrambe le letture.
    \item Come membro del gruppo, è la riscrittura del tutorial della slide precedente: $\lambda_i=0$ e $u_i=x_i(0)$. Come opinione fissa esterna, è il pregiudizio degli agenti che la ascoltano, come qui.
    \item Le due letture danno le stesse traiettorie. E il disaccordo dei due agenti stubborn è il disaccordo di agenti ancorati a pregiudizi diversi.
    \item Ora gli esempi del modello FJ, sulle stesse due reti.
  \end{discorso}
\end{slide}

% =====================================================================
\parte{Parte 4 — Il modello FJ, esempi}{5:00}

\begin{slide}{m-one-prej}{A single prejudice}{0:50}
  \begin{discorso}
    \item Stesso peso un mezzo sull'agente 1 nei due modelli: a sinistra come self-weight, nel FDG con $A(\tfrac12)$; a destra come attaccamento al pregiudizio, nel FJ con $\lambda_1=\tfrac12$ e $u=x(0)$.
    \item La rete è fortemente connessa, quindi tutti gli agenti sono P-dependent e il modello FJ è asintoticamente stabile, anche se $G$ è periodico e non ha self-loop.
    \item Nel FJ $V=\uno e_1^\top$: tutto il gruppo finisce sul pregiudizio dell'agente 1, zero, per ogni attaccamento positivo, per quanto piccolo. Nel FDG invece il consenso è a $3/7$.
    \item Sulla misura del potere: $c=e_1$ per ogni $\lambda_1<1$, l'agente 1 prende tutta la influence centrality; nel FDG, a $t=\tfrac12$, la sua social power è $2/7$.
    \item È di nuovo la testardaggine che ribalta la dinamica: un solo agente ancorato, che influenza tutti, decide dove finisce tutto il gruppo, per quanto debole sia il suo attaccamento.
    \item E se tutti gli agenti hanno un pregiudizio?
  \end{discorso}
\end{slide}

\begin{slide}{m-all-prej}{Every agent prejudiced}{0:45}
  \begin{discorso}
    \item Questo è il confronto equo: lo stesso peso 0.2 per tutti gli agenti. A sinistra il FDG con $A=0.2I+0.8W$, a destra il FJ con $\Lambda=0.8I$ e $u=x(0)$.
    \item Nel FDG la matrice è primitiva, per il Corollario 1.12, e doubly-stochastic: consenso sulla media, 0.5.
    \item Nel FJ tutti gli agenti sono prejudiced, quindi il modello è asintoticamente stabile, e $V=0.2(I-0.8W)^{-1}$ ha tutte le entrate positive. I pregiudizi sono diversi, quindi disaccordo.
    \item Su una rete fortemente connessa il FDG non va mai in disaccordo: nel FJ il disaccordo viene dai pregiudizi.
    \item In termini sociali: l'inerzia alla fine cede, l'ancora no; ancore diverse tengono il gruppo diviso, quello che il tutorial chiama un cleavage di opinioni.
    \item La slide seguente mostra perché questo è un disaccordo anche nel senso del corso.
  \end{discorso}
\end{slide}

\begin{slidedue}{m-six-sinks}{Six prejudices, six sinks}{1:15}
  \begin{discorso}
    \item Riscrivo il modello FJ della Figura 10 come modello FDG su dodici agenti. Ogni agente reale dà peso totale 0.8 agli altri agenti reali e 0.2 al suo agente virtuale, che è stubborn sul pregiudizio $u_i$.
    \item Gli agenti reali formano una sola componente fortemente connessa, quella ombreggiata, che non è un sink. Collassandola si ottiene il condensation digraph: i sei agenti virtuali sono sei sink, aperiodici per i self-loop.
    \item Teorema 2.4 con sei sink: è il caso (iv), come nel Remark 2.7. Lo stesso vale sull'anello.
    \item Quindi $x(\infty)=Vu$, con $V$ a entrate positive: ogni opinione finale è una combinazione convessa di tutti e sei i pregiudizi.
    \item Anche la Figura 9 ha sei sink, ma solo $1'$ è ascoltato: togliendo gli altri cinque resta un solo sink, e il consenso.
    \item Così si chiude il filo del disaccordo del FDG: un agente stubborn dà il consenso, nella Figura 2 con $t=1$; due danno il disaccordo, l'esempio della domanda; sei agenti virtuali, tutti ascoltati, danno il disaccordo di questa slide.
    \item \clic Una nota sulle traiettorie. La Figura 10 mostra solo i sei agenti reali; nel modello aumentato le traiettorie sarebbero dodici, e quelle dei sei agenti virtuali sono rette orizzontali sui pregiudizi.
    \item Ogni agente reale parte sulla retta del suo agente virtuale e la lascia appena la media delle opinioni che ascolta è diversa: sul ciclo subito, sull'anello gli agenti da 2 a 5 più tardi, quando arriva il cambiamento degli agenti 1 e 6.
    \item Il peso 0.2 sul suo agente virtuale lo richiama sempre, e si ferma su $Vu$, mai sul proprio pregiudizio. Nel FDG lo stesso 0.2 sta sull'opinione attuale, e l'agente segue il gruppo fino al consenso.
    \item Ultimo caso: dove nessun pregiudizio arriva.
  \end{discorso}
\end{slidedue}

\begin{slidedue}{m-aux}{Where no prejudice arrives}{1:10}
  \begin{discorso}
    \item Questa rete serve per il caso escluso finora: non tutti gli agenti sono P-dependent.
    \item Gli agenti 1, 2 e 3 formano un gruppo chiuso, un sink del condensation digraph; 4, 5 e 6 sono una catena che lo ascolta. Un solo agente è prejudiced, con $\lambda=0.5$; ombreggiati gli agenti P-dependent.
    \item Con $W_p$ il gruppo chiuso è un ciclo periodico; con $W_a$ l'arco da 2 a 1 lo rende aperiodico.
    \item \clic Caso (a): se il prejudiced è l'agente 2, dentro il gruppo chiuso, tutti sono P-dependent. Il modello è asintoticamente stabile e tutti finiscono su $u_2=0.5$, anche se il ciclo è periodico.
    \item Caso (b): se il prejudiced è l'agente 5, fuori, solo 4 e 5 sono P-dependent. Il gruppo chiuso è periodico, $W^{22}$ non è regular, e il modello non converge, per il Teorema 3.2.
    \item Caso (c): con l'arco da 2 a 1, $W^{22}$ è regular. Il modello converge, ma non è asintoticamente stabile: gli agenti 1, 2, 3 e 6 finiscono a 0.4, gli agenti 4 e 5 a 0.65.
    \item Caso (d): se cambio solo $x_2(0)$, da 0.5 a 1, i limiti diventano 0.6 e 0.75. Il limite dipende dalle opinioni iniziali dei P-independent.
    \item In termini sociali: gli agenti che nessun pregiudizio raggiunge formano un modello FDG a sé; hanno bisogno dell'aperiodicità e ricordano le loro opinioni iniziali.
    \item Infine, le due misure di potere.
  \end{discorso}
\end{slidedue}

\begin{slide}{m-lemma}{Centrality meets social power}{1:00}
  \begin{discorso}
    \item Per confrontare le due misure prendo $W=A(\tfrac12)$, la rete con self-weight un mezzo sull'agente 1, e la stessa suscettibilità $\alpha$ per tutti, con $u=x(0)$.
    \item $W$ è primitiva, per il Corollario 1.12, quindi fully regular: il Lemma 3.5 si applica.
    \item Per $\alpha=0$, $V=I$ e la influence centrality è uniforme, $1/6$. Al crescere di $\alpha$, $c_1$ sale fino alla social power $w_1=2/7$, e $c_2$ scende fino a $w_2=1/7$. Sul ciclo $c_2$ scende prima anche sotto $w_2$.
    \item In termini sociali: man mano che gli agenti si aprono agli altri, l'inerzia dell'agente 1 conta di più.
    \item La social power dice chi guida un gruppo che si mette d'accordo; la influence centrality dice chi pesa di più in un gruppo che resta in disaccordo. La seconda diventa la prima quando le ancore si allentano.
    \item Per concludere, i due modelli a confronto.
  \end{discorso}
\end{slide}

% =====================================================================
\parte{Parte 5 — Conclusioni}{1:30}

\begin{slide}{m-table}{FDG and FJ side by side}{0:40}
  \begin{discorso}
    \item La tabella mette a confronto i due modelli riga per riga.
    \item La resistenza: nel FDG è un self-weight, un'inerzia; nel FJ un attaccamento, un'ancora.
    \item L'agente stubborn: nel FDG è una condizione sulla sola matrice, $a_{ii}=1$; nel FJ sulla matrice e sull'ingresso, $\lambda_i=0$ e $u_i=x_i(0)$. La sua opinione è uno stato nel FDG, un ingresso nel FJ.
    \item La stabilità asintotica: mai nel FDG; nel FJ se e solo se tutti gli agenti sono P-dependent.
    \item Su una rete fortemente connessa il FDG dà consenso o nessuna convergenza, mai disaccordo; il FJ dà disaccordo ogni volta che i pregiudizi che contano sono diversi.
    \item L'aperiodicità serve sempre nel FDG; nel FJ solo per gli agenti P-independent.
    \item E la misura del potere: la social power richiede il consenso; la influence centrality richiede la stabilità asintotica e $u=x(0)$.
  \end{discorso}
\end{slide}

\begin{slide}{m-messages}{Take-home messages}{0:50}
  \begin{discorso}
    \item Tre messaggi da portare a casa.
    \item Primo, inerzia e ancora. Un self-weight rallenta un agente ma lo lascia seguire gli altri: su una rete fortemente connessa e periodica basta un self-weight fra 0 e 1 per trasformare le oscillazioni in consenso. Un pregiudizio non cede mai: un solo agente prejudiced che influenza tutti decide dove finisce il gruppo, e pregiudizi diversi lo tengono diviso.
    \item Secondo, stato e ingresso. La stessa opinione stubborn è uno stato nel modello FDG, che non è mai asintoticamente stabile, e un ingresso nel modello FJ, asintoticamente stabile quando tutti gli agenti sono P-dependent. Il grafo aumentato riporta gli ingressi a stati.
    \item Terzo, il potere. La social power cresce gradualmente con il self-weight; un agente prejudiced che influenza tutti prende tutta la influence centrality; e con $\Lambda=\alpha I$ e $W$ fully regular, la influence centrality tende alla social power quando $\alpha$ tende a 1.
  \end{discorso}
\end{slide}

\begin{slidepag}{\the\numexpr\pg{m-messages}+1\relax}{Thank you!}{}
  \begin{discorso}
    \item Grazie per l'attenzione. Sono a disposizione per le domande.
  \end{discorso}
\end{slidepag}

\end{document}
